% Figure: Mitochondrial electron transport chain + ATP synthase,
% with proton-motive-force diagram.
\begin{figure}[htbp]
\centering
\begin{tikzpicture}[
  font=\footnotesize,
  complex/.style={draw=engblack!60, fill=engblue!15,
    minimum width=1.8cm, minimum height=2.2cm,
    rounded corners=2pt, align=center, inner sep=2pt},
  syn/.style={draw=engblack!60, fill=englime!25,
    minimum width=1.6cm, minimum height=2.2cm,
    rounded corners=2pt, align=center, inner sep=2pt},
  arrow/.style={-{Stealth[length=2mm]}, draw=engamber, thick},
  flow/.style={-{Stealth[length=1.6mm]}, draw=engblue, dashed},
]
% Inner membrane: two horizontal lines
\draw[engblack!50, line width=1pt] (-0.5,1.1) -- (12.5,1.1);
\draw[engblack!50, line width=1pt] (-0.5,-1.1) -- (12.5,-1.1);
% Membrane shading
\fill[engblack!8] (-0.5,-1.1) rectangle (12.5,1.1);

\node[align=left, font=\scriptsize\itshape, color=engblack!60]
  at (-0.5,1.6) {Intermembrane space\\ (high [H$^+$])};
\node[align=left, font=\scriptsize\itshape, color=engblack!60]
  at (-0.5,-1.6) {Matrix\\ (low [H$^+$])};

% Four complexes + ATP synthase
\node[complex] (c1) at (1,0) {\textbf{I}\\NADH\\dehydrog.};
\node[complex] (c2) at (3,0) {\textbf{II}\\Succinate\\DH};
\node[complex] (c3) at (5.5,0) {\textbf{III}\\Cyt $bc_1$};
\node[complex] (c4) at (8,0) {\textbf{IV}\\Cyt $c$\\oxidase};
\node[syn]     (atps) at (11,0) {\textbf{ATP}\\synthase};

% Electron flow (dashed blue arrows)
\draw[flow] (c1.east) -- (c2.west);
\draw[flow] (c2.east) -- (c3.west);
\draw[flow] (c3.east) -- (c4.west);

% Substrate / product labels
\node[align=center, font=\scriptsize] at (1, -1.8) {NADH $\to$\\NAD$^+ + $ H$^+$};
\node[align=center, font=\scriptsize] at (3, -1.8) {FADH$_2 \to$\\FAD};
\node[align=center, font=\scriptsize] at (8, -1.8) {$\tfrac{1}{2}$O$_2 + 2$H$^+$\\$\to$ H$_2$O};
\node[align=center, font=\scriptsize] at (11, -1.8) {ADP $+$ P$_i$\\$\to$ ATP};

% Proton pumping arrows from matrix to IMS at I, III, IV
\draw[arrow] (1,-1.0) -- (1,1.05);
\draw[arrow] (5.5,-1.0) -- (5.5,1.05);
\draw[arrow] (8,-1.0) -- (8,1.05);
\node[font=\scriptsize, color=engamber] at (1.4,0.0) {4 H$^+$};
\node[font=\scriptsize, color=engamber] at (5.9,0.0) {4 H$^+$};
\node[font=\scriptsize, color=engamber] at (8.4,0.0) {2 H$^+$};

% H+ flow through ATP synthase from IMS to matrix
\draw[arrow, draw=engred] (11,1.05) -- (11,-1.0);
\node[font=\scriptsize, color=engred] at (11.5,0.0) {$\sim$3 H$^+$};

% PMF caption box, below the membrane
\node[draw=engblack!60, fill=engblue!8, rounded corners=2pt,
  align=left, font=\scriptsize, text width=11cm, inner sep=4pt]
  at (5.8,-3.2)
  {\textbf{Proton-motive force:} $\Delta p = \Delta \Psi - (2.303 RT / F)\,\Delta\text{pH}
   \approx -200$~mV (matrix negative, alkaline). Free energy per
   proton: $\Delta G_{H^+} = F \Delta p \approx -19$~kJ/mol. ATP
   synthase couples $\sim 3$ H$^+$ in to one ATP synthesised:
   $\Delta G_{\text{ATP}} = -3 \cdot \Delta G_{H^+} \approx +57$~kJ/mol
   (matches the cellular ATP free energy of $\sim 55$~kJ/mol).};
\end{tikzpicture}
\caption[Electron transport chain and chemiosmotic coupling]{The
mitochondrial inner membrane with the four respiratory complexes (I,
II, III, IV) and ATP synthase. Electrons (dashed blue) enter at
Complex I (from NADH) or Complex II (from FADH$_2$), traverse the
chain, and reduce O$_2$ to water at Complex IV. Complexes I, III, and
IV pump protons (amber, $4+4+2 = 10$ protons per NADH) from the
matrix to the intermembrane space, building the proton-motive force.
Protons return through ATP synthase (red), turning the rotor that
condenses ADP and P$_i$ into ATP. The stoichiometry of $\sim 3$
protons per ATP gives the cellular yield of $\sim 2.5$ ATP per NADH
(matching $10/3 \approx 3.3$ protons-per-NADH-to-ATP). FADH$_2$ enters
at Complex II, which does not pump, and yields $\sim 1.5$ ATP. After
Mitchell \cite{paper:mitchell1961-chemiosmotic} and Nelson and Cox
\cite[ch.~19]{text:lehninger-biochemistry}.}
\label{fig:vol06:ch02:etc}
\end{figure}
